\section{Discussion and extensions}\label{sec:discussion}

The household stage reports three things the first pass could only
frame: the modelled cash benefits do not respond to a price shock
while indexation ran a year late (Section~\ref{sec:results-rulebook});
the discretionary instruments cushioned the energy shock at rates
that fall with income, while the TCA food shock, with a similar
spending-share gradient, had no instrument legislated against it
(Section~\ref{sec:results-tca-ss}; the two episodes differed in size
and speed, and the paper does not estimate why the responses
differed); and most of the variance in the energy burden lies
\emph{within} income deciles, the pattern from which
\citet{levell2025} conclude that an income-conditioned instrument
targets energy exposure poorly.

Three descriptive findings hold across the paper's variants. The
negative shocks of the 2020s, two policy-made, one
trade-transmitted, were regressive through prices, and the projected
gains from the positive policy episode are similar shares of spending
across deciles and smaller in aggregate, because they arrive in
categories (clothing, footwear) whose budget shares are flat across
the distribution. The
modelled responses ran through different mechanisms: the automatic
tax--benefit system where the shock hit earnings (in the declared
tariff scenario), discretionary instruments where it hit prices. The
largest discretionary instrument, capping a price, cushioned
proportionally by construction, and the full modelled stack on
household records cushioned at rates that fall with income
(Section~\ref{sec:results-secondstage}). The welfare literature reads
proportional energy support two ways: \citet{levell2025} retain a
positive price subsidy in their optimal package because income
predicts energy exposure poorly, while \citet{bonnet2025} find that
transfers dominate subsidies for transport fuel, whose budget shares
rise with income; the eight UK energy support schemes cost \pounds 44
billion against an initial \pounds 139 billion estimate, of which the
Guarantee was \pounds 24.0 billion \citep{nao2024}. A domestic
incidence exercise cannot price the externality that national
subsidies impose through world energy demand \citep{auclert2023}.
Benefit indexation, the automatic price-side channel, followed the
previous September's inflation, so its April 2022 step fell below the
inflation of the year in which the shock landed.

Past shocks carry ex-post estimates and recent agreements ex-ante
projections. The sources reviewed for this paper include no published
UK household-level ex-post validation of a trade agreement's impact
assessment; the Department for Business and Trade's CPTPP assessment
\citep{dbtcptpp2023}
schedules an evaluation within five years of entry into force, and the
Regulatory Policy Committee asked for the India assessment's consumer
analysis to be developed further \citep{rpc2025}. On the government's
own projections, the positive policy episode's gains are small and
distributionally flat; on the published estimates, the negative
episodes are large and regressive. A household-level ex-post analysis
would put both on a comparable basis.

The limits are those declared throughout. First stages are imported,
with their identification and their controversies; the incidence is
first-order and static; the channels are not summed; the rulebook
decomposition labels legislated change as discretion without asking
what would have been legislated anyway. Three price-led episodes are carried to household records and the
tariff episode's earnings channel is computed in this repository; the deeper
rulebook and poverty outputs remain, and the first-pass numbers
should be read accordingly. Poverty
and inequality effects, the remaining episodes' rulebook
decompositions, and the consumption channel's interaction with the
earnings channel on common households all remain for the extension
Section~\ref{sec:roadmap} specifies.

The framework can be reused: a future trade shock that fits its price
or earnings channels enters as one more declared first stage for the
same second stage, subject to the data and model limits stated above.
